\subsection{谱的孤立点与敏感点}

设 E 是完备可赋范空间，u 是 E 的一个自同态。设 \(\lambda\) 是 u 的谱的一个孤立点。回忆我们用 \(e_{\lambda}(u)\) 表示与 u 及 u 的谱的既闭又开的子集 \(\{\lambda\}\) 相伴的谱投影（I, p. 131 的第 3 目）。自同态 \(u - \lambda1_{E}\) 诱导 \(\operatorname{Ker}(e_{\lambda}(u))\) 的一个自同构以及 \(\operatorname{Im}(e_{\lambda}(u))\) 的一个拟幂零自同态（同前）。

定义 1. --- 我们说 \(\lambda\) 对 \(u\) 有有限谱重数，如果谱投影 \(e_{\lambda}(u)\) 有限秩；此时，整数 \(m_{\lambda}(u) = \dim (\mathrm{Im}(e_{\lambda}(u)))\) 称为 \(\lambda\) 对 \(u\) 的谱重数。

定义 2. --- 我们把 Sp(u) 的有限谱重数的孤立点称为 u 的谱的敏感点。Sp(u) 的敏感点所成的集称为 u 的敏感谱，记作 Sp \(_{s}\) (u)。\(^{(2)}\) 我们把 u 的本质谱记作 Sp \(_{e}\) (u)，即补集 \(\mathrm{Sp}(u) - \mathrm{Sp}_{\mathrm{s}}(u)\)。

由于集 \(\mathrm{Sp}_{\mathrm{s}}(u)\) 由 \(\mathrm{Sp}(u)\) 的孤立点组成，它在 \(\mathrm{Sp}(u)\) 中是开的、离散的且可数的（III, p. 78, 引理 3）；它还是有界的。集 \(\mathrm{Sp}_{\mathrm{e}}(u)\) 在 \(\mathrm{Sp}(u)\) 中是闭的，故是紧的。

由 III, p. 54 的命题 14，\(u\) 的谱的敏感点是那些复数 \(\lambda\) （使 \(u - \lambda 1_{\mathrm{E}}\) 为 E 的里斯自同态且不是 E 的自同构的）。

对每个复数 \(\lambda\) ，用 \(\mathrm{N}_{\lambda}(u)\) 与 \(\mathrm{I}_{\lambda}(u)\) 表示 E 的自同态 \(u - \lambda 1_{\mathrm{E}}\) 的幂零空间与余幂零空间。由于 E 是弗雷歇空间，\(u - \lambda 1_{\mathrm{E}}\) 是里斯自同态当且仅当 \(\mathrm{N}_{\lambda}(u)\) 有限维且 \(\mathrm{I}_{\lambda}(u)\) 有有限余维数（III, p. 53 的命题 13）；在此情形，\(u - \lambda 1_{\mathrm{E}}\) 可逆当且仅当 \(\mathrm{N}_{\lambda}(u)\) 缩减为 0（III, p. 47, 命题 6）。于是 \(u\) 的谱的敏感点是复数 \(\lambda\) （使 \(\mathrm{N}_{\lambda}(u)\) 为非零有限维且 \(\mathrm{I}_{\lambda}(u)\) 在 E 中有有限余维数的）（III, p. 54, 命题 14）。此时谱重数 \(m_{\lambda}(u)\) （\(\lambda\) 的）就是 \(\mathrm{N}_{\lambda}(u)\) 的维数，且自同态 \(u - \lambda 1_{\mathrm{E}}\) 通过限制诱导 \(\mathrm{I}_{\lambda}(u)\) 的一个自同构与 \(\mathrm{N}_{\lambda}(u)\) 的一个幂零自同态。由于 \(\mathrm{N}_{\lambda}(u)\) 不缩减为 0，由此得出 \(\lambda\) 是 \(u\) 的一个特征值。最小的整数 \(n \geqslant 1\) （使 \(\operatorname{Ker}((u - \lambda 1_{\mathrm{E}})^n)\) 等于 \(\mathrm{N}_{\lambda}(u)\) ）是 \(u\) 的预解式在点 \(\lambda\) 处极点的阶（I, p. 131 的第 3 目），它以 \(m_{\lambda}(u)\) 为上界。

特别地，我们由此证明了：

命题 1. --- u 的谱的敏感点是有限谱重数的特征值。

\(u\) 的本质谱由那些复数 \(\lambda\) （使 \(u - \lambda 1_{\mathrm{E}}\) 不是 E 的里斯自同态的）组成。特别地，若 \(u\) 是紧的，则每个 \(\lambda \in \mathrm{Sp}(u) - \{0\}\) 都是谱的敏感点（III, p. 77 的推论 2）。

命题 2. --- 设 E 是完备可赋范空间，u 是 E 的一个自同态。设 H 是 \(\mathrm{Sp}_{\mathrm{s}}(u)\) 的一个有限子集。集 H 在 \(\mathrm{Sp}(u)\) 中既开又闭。谱投影 \(e_{\mathrm{H}}(u)\) 有限秩，其像为 \(\sum_{\lambda \in \mathrm{H}}\mathrm{N}_{\lambda}(u)\) ，核为 \(\bigcap_{\lambda \in \mathrm{H}}\mathrm{I}_{\lambda}(u)\) 。

这由 I, p. 131 的第 3 目得出，因为对每个 \(\lambda \in \mathrm{Sp}_{\mathrm{s}}(u)\) ，子空间 \(\mathrm{N}_{\lambda}(u)\) 与 \(\mathrm{I}_{\lambda}(u)\) 分别是谱投影 \(e_{\lambda}(u)\) 的像与核。

推论 1. --- 设 V 是 \(\mathrm{Sp_e}(u)\) 的一个邻域。

\begin{enumerate}
\def\labelenumi{\alph{enumi})}
\item
  存在 E 作为拓扑直和 F ⊕ G 的分解，使得 F 有限维，F 与 G 在每个与 u 交换的自同态下稳定，且 u 诱导的 G 的自同态的谱包含于 V。
\item
  假设 V 非空。存在元素 v （属于 u 在 \(\mathcal{L}(\mathrm{E})\) 中的双交换子），其谱包含于 V 且 \(v - u\) 有限秩。
\end{enumerate}

令 \(\mathrm{H} = \mathrm{Sp}(u) \cap (\mathbf{C} - \mathring{\mathrm{V}})\) 。集 H 包含于 \(\mathrm{Sp}_{\mathrm{s}}(u)\) ，故是离散的；由于它在紧空间 \(\mathrm{Sp}(u)\) 中是闭的，它是有限的。可取谱投影 \(e_{\mathrm{H}}(u)\) 的像与核作为 F 与 G（命题 2）。

假设 V 非空。设 \(\mu\in\mathrm{V}\) ，用 \(v\) 表示 E 的这样一个自同态：它在 F 上与比率为 \(\mu\) 的位似重合，在 G 上与 \(u\) 重合。其谱包含于 \(\{\mu\}\cup(\mathrm{Sp}(u)-\mathrm{H})\) ，从而包含于 V，且 \(v-u\) 有限秩，因为 F 有限维。

对 C 的每个紧子集 S，我们用 \(\mathcal{O}(S)\) 表示在 S 的一个邻域中定义、取值于 C 的全纯函数芽所成的代数（I, p. 49, §4, 第 1 目）。

推论 2. --- 设 \(f \in \mathcal{O}(\mathrm{Sp}(u))\) ，\(\mu\) 是一个复数。设 H 是 \(\lambda \in \mathrm{Sp}(u)\) （满足 \(f(\lambda) = \mu\) ）所成的集。则 \(\mu\) 是 \(f(u)\) 的谱的敏感点，当且仅当 H 有限、非空且包含于 \(u\) 的敏感谱。在这些条件下，谱投影 \(e_{\mu}(f(u))\) 等于投影 \(e_{\mathrm{H}}(u)\) （与 \(u\) 及 H 相伴的）。我们有

\[
\mathrm{N} _ {\mu} (f (u)) = \bigoplus_ {\lambda \in \mathrm{H}} \mathrm{N} _ {\lambda} (u), \quad \mathrm{I} _ {\mu} (f (u)) = \bigcap_ {\lambda \in \mathrm{H}} \mathrm{I} _ {\lambda} (u)
\]

而 \(\mu\) 对 \(f(u)\) 的谱重数是 H 的元素的谱重数之和，即

\[
m _ {\mu} (f (u)) = \sum_ {\lambda \in \mathrm{H}} m _ {\lambda} (u).
\]

复数 \(\mu\) 属于 \(f(\mathrm{Sp}(u))\) ，从而属于 \(f(u)\) 的谱（I, p. 75 的命题 8），当且仅当 H 非空。此时我们有 \(e_{\mu}(f(u)) = e_{\mathrm{H}}(u)\) （I, p. 127 的第 1 目）。其余断言由命题 2 得出。

例. --- 1) 设 E 是有限维向量空间，u 是 E 的一个自同态。u 的谱有限且与 \(\mathrm{Sp}_{\mathrm{s}}(u)\) 重合。其元素是 u 的特征多项式 \(\chi_{u}\) 的根；它们就是 u 的特征值。谱重数 \(m_{\lambda}(u)\) （一个元素 \(\lambda \in \mathrm{Sp}(u)\) 的）是 \(\lambda\) 作为多项式 \(\chi_{u}\) 的根的重数（A, VII, p. 36, 推论）。由上面的推论 2，我们有

\[
\chi_ {f (u)} (\mathrm{T}) = \prod_ {\lambda \in \operatorname{Sp} (u)} (\mathrm{T} - f (\lambda)) ^ {m _ {\lambda}}
\]

（对一切 \(f \in \mathcal{O}(\mathrm{Sp}(u))\) ）。这推广了 A, VII, p. 36 的命题 10。

\begin{enumerate}
\def\labelenumi{\arabic{enumi})}
\setcounter{enumi}{1}
\tightlist
\item
  设 X 是 C 的一个紧子集，\(\mathcal{C}(X)\) 是 X 上取复值的连续函数所成的巴拿赫空间。用 u 表示 \(\mathcal{C}(X)\) 的这样一个自同态：它把 \(f \in \mathcal{C}(X)\) 映为函数 \(z \mapsto zf(z)\) （在 \(\mathcal{C}(X)\) 中） 。u 的谱等于 X，其敏感点是 X 的孤立点，且它们的谱重数等于 1。
\end{enumerate}

